For every finite rational-coefficient input \((G,\mathcal A,V,W,T,Y)\), the corrected guarded-district handle defines a terminating four-output fixed-alphabet front end \(\mathsf{F\text{-}ID}^{+}\). It first returns the feasibility or infeasibility certificate of Theorem \(\mathrm{thm:feasibility\text{-}certificate}\). Conditional on feasibility, it enumerates functional-map and row-support strata, compiles their observable supports into guards in \(p\), and does exactly one of the following: (i) if every feasible observable-support branch is covered by a coefficient-verified conditional-event-ratio formula or by the one-way classical-\(\mathsf{ID}\) embedding, it returns a support-guarded rational certificate \(C\) for the complete intervention vector \(q_M\) for every \(M\in\mathfrak M(G,\mathcal A,V,W,T)\); (ii) if an observationally unreachable parent row is target-relevant, it returns an original-alphabet rational separating certificate \(\mathsf{Sep}\); or (iii) it returns the already-certified success leaves together with an exact residual positive-support semialgebraic fiber formula and makes no identification verdict on that formula. More precisely, after target-inert unreachable rows have been fixed canonically and \(\mathcal U\) denotes the remaining unresolved strata, the residual output is
\[
 \Psi_{\mathcal U}(p)
 :=
 \bigvee_{s_0,s_1\in\mathcal U}
 \exists\theta^{(0)},\theta^{(1)}
 \left\{
 \begin{aligned}
 &\overline\Phi_{s_0}(\theta^{(0)})\wedge
   \overline\Phi_{s_1}(\theta^{(1)})\\
 &{}\wedge\bigwedge_{v\in\mathcal A_V}
   \bigl(P_{\theta^{(0)}}(v)=p(v)=P_{\theta^{(1)}}(v)\bigr)\\
 &{}\wedge\bigvee_{t\in\mathcal A_T,\,y\in\mathcal A_Y}
   \bigl(Q_{\theta^{(0)}}(t,y)\ne Q_{\theta^{(1)}}(t,y)\bigr).
 \end{aligned}
 \right.
\]
Here every free, noncanonical row in \(\overline\Phi_s\) has a strictly positive observational parent-configuration probability. The formula \(\Psi_{\mathcal U}(p)\) is true exactly when the unresolved part of the fiber at \(p\) contains two models with different targets. The procedure neither decides \(\Psi_{\mathcal U}\) nor invokes generic real quantifier elimination to search for an identification verdict; if satisfiability is established independently, the real-algebraic separator lemma converts it to an exactly encodable algebraic pair.

Let \(B=F2^K\). A direct implementation uses
\[
 \begin{aligned}
 \mathcal T={}&O\!\left(FJ_T(d+m)h+FR+4^d(d+m+h)\right)\\
 &{}+O\!\left(BJ_T(d+m)\left[h^3 4^d6^h+Rh^4(d+1)^K\right]\right)\\
 &{}+O\!\left(B^2\left[2R+4K+2J_T+(2h+2J_Th)dh\right]\right)
 \end{aligned}
\]
table and exact rational-arithmetic operations and has output length \(O(\mathcal T)\) rational words. Since \(R\le K\le dh\), this is an explicit elementary bound in \(F,J_T,2^K,h,d,m\). In the expanded residual formula there are at most \(2K\le2dh\) existential CPT variables, at most \(h\) free observed-law coordinates, at most \(B^2\) stratum-pair disjuncts, and at most
\[
 B^2\bigl(2R+4K+2J_T+2h+2J_Th\bigr)
\]
polynomial sign atoms. Their degrees are at most \(d\), and their integer coefficient bit heights are at most \(\lceil\log_2(2h+1)\rceil\). A support-hole pair constructed by the scan has CPT denominators dividing \(L=h(d+1)+1\), so its CPT-entry bit height is at most \(\lceil\log_2 L\rceil+1\), while direct exact checks of its observed and intervention polynomials have bit length \(O(d\log L+\log h)\). Thus this theorem proves a total exact four-output front end, not the stronger universal success-or-countermodel dichotomy.